% Options for packages loaded elsewhere
\PassOptionsToPackage{unicode}{hyperref}
\PassOptionsToPackage{hyphens}{url}
%
\documentclass[11pt]{article}
\usepackage{amsmath,amssymb}
\usepackage{iftex}
\ifPDFTeX
  \usepackage[T1]{fontenc}
  \usepackage[utf8]{inputenc}
  \usepackage{textcomp}
\else
  \usepackage{unicode-math}
  \defaultfontfeatures{Scale=MatchLowercase}
  \defaultfontfeatures[\rmfamily]{Ligatures=TeX,Scale=1}
\fi
\usepackage{lmodern}
\usepackage{xcolor}
\usepackage{booktabs}
\usepackage{array}
\usepackage{graphicx}
\usepackage{float}
\usepackage{geometry}
\geometry{a4paper, margin=1in}
\usepackage{bookmark}
\hypersetup{
  hidelinks,
  pdfcreator={LaTeX}
}
\author{}
\date{}

\begin{document}

\section{Part 3: Warehouse Operations Solutions}\label{part3}

\subsection{Question 3.1: Slotting Optimization}\label{q3.1}

Using the ABC-XYZ classification from Part 1 and the physical SKU attributes from the
Master Data, we assigned all 470 classified SKUs to three warehouse zones.
We present two models: \textbf{Model 1} (ABC-only baseline) and
\textbf{Model 2} (ABC + Ergonomics + Pick-Mix), which achieves the
\textbf{30\% peak-time improvement target} through ergonomic sub-tiering.

\subsubsection{Model 1 --- ABC-Only Baseline}\label{model1}

\begin{itemize}
    \item \textbf{Pick-Face Zone (Ground Level)} --- Class A SKUs: lowest racks at
          ergonomic height, nearest to the I/O dock and packing station.
    \item \textbf{Forward Reserve Zone} --- Class B SKUs: two-bin replenishment logic.
          A forward pick-face bin is replenished from the bulk reserve (WMS-triggered).
    \item \textbf{Reserve / Bulk Zone} --- Class C SKUs: upper racks and pallet block-stack.
\end{itemize}

\begin{table}[H]
\centering
\caption{Model 1 --- Zone Summary (ABC-Only)}
\begin{tabular}{lcccc}
\toprule
Zone & Class & SKU Count & \% of Picks & \% of Volume \\
\midrule
        Pick-Face Zone & A & 28 & 27.7\% & 69.83\% \\
        Forward Reserve Zone & B & 39 & 38.2\% & 20.04\% \\
        Reserve / Bulk Zone & C & 403 & 34.1\% & 10.13\% \\
\bottomrule
\end{tabular}
\end{table}

The Pick-Face Zone contains 28 SKUs (6.0\% of the assortment) but generates
\textbf{27.7\%} of all pick transactions and
\textbf{69.8\%} of total volume.
Within this zone, the \textbf{19 Class A-X fast-movers} occupy the first 19 rank slots
and drive \textbf{25.88\%} of pick frequency and \textbf{57.87\%} of volume.

\subsubsection{Model 2 --- ABC + Ergonomics + Pick-Mix}\label{model2}

\paragraph{Observed pick-unit mix.}
Analysis of the transaction log (Jun--Dec 2025) reveals that Class A pick orders are
\textbf{predominantly carton picks (58\%)} rather than pallet picks, and the two sub-populations
have fundamentally different handling requirements:

\begin{table}[H]
\centering
\caption{Pick-Unit Mix by ABC Zone (avg. order-line share)}
\begin{tabular}{lrrr}
\toprule
Zone & Pallet pick & Carton pick & Loose pick \\
\midrule
        Pick-Face Zone (A) & 36\% & 58\% & 6\% \\
        Forward Reserve (B) & 14\% & 61\% & 25\% \\
        Reserve / Bulk (C) & 2\% & 37\% & 61\% \\
\bottomrule
\end{tabular}
\end{table}

\paragraph{Sub-tier logic.}
Model 2 splits each ABC zone into sub-tiers based on pick-unit mix and physical dimensions:

\begin{table}[H]
\centering
\caption{Model 2 --- Sub-Tier Assignment Rules and SKU Counts}
\begin{tabular}{llp{6cm}r}
\toprule
Sub-Tier & Name & Assignment Rule & SKUs \\
\midrule
        A1 & Pallet Lane & pallet\_share $\geq$ 40\% AND CBM $\geq$ 0.07\,m$^3$ --- ground-level forklift slots & 3 \\
        A2 & Big-Face Shelf & pallet\_share $<$ 40\% AND carton\_share $\geq$ 80\% --- waist-height shelving & 7 \\
        A3 & Mixed-Pick & remaining Class A & 18 \\
        \midrule
        B1 & Bulk-Replen & pallet\_share $\geq$ 30\% OR CBM $\geq$ 0.07\,m$^3$ --- mid-aisle ground storage & 9 \\
        B2 & Two-Bin & remaining Class B --- standard two-bin forward pick-face & 30 \\
        \midrule
        C1 & Upper Rack & CBM $<$ 0.05\,m$^3$ --- small/light slow-movers & 339 \\
        C2 & Pallet Block & CBM $\geq$ 0.05\,m$^3$ --- floor block-stack behind Zone B & 64 \\
\bottomrule
\end{tabular}
\end{table}

\paragraph{Composite intra-zone score.}
Within each sub-tier, slots are ranked by a composite score:
\[
score = 0.50 \times v_{rank} + 0.30 \times ergo_{penalty} + 0.20 \times size_{score}
\]
where $v_{rank}$ is the normalised velocity rank (0=fastest), $ergo_{penalty}$
is 0 / 0.5 / 1.0 for pcs weight $\leq$3\,kg / $\leq$8\,kg / $>$8\,kg, and
$size_{score}$ is CBM relative to the median within the same ABC class.

\paragraph{A-Door Rule.}
The top-10 velocity SKUs (by order frequency) are pinned to the first 10 rack bays
from the outbound dock regardless of sub-tier --- eliminating travel for the most
frequently picked items during peak-hour bursts.

\subsubsection{Top Pick-Face SKUs (Class A, by Frequency)}\label{top-a}

\begin{table}[H]
\centering
\caption{Top 5 Class-A SKUs assigned to Pick-Face Zone}
\begin{tabular}{clllrr}
\toprule
Rank & SKU Code & Product & Class & Picks (6M) & Volume (units) \\
\midrule
        1 & 4200008229 & Quạt treo ASIAvina - VY377990 - Xám & A-X & 825 & 5,608 \\
        2 & 4200008371 & QUẠT ĐỨNG VY619990 & A-X & 682 & 3,355 \\
        3 & 4200008366 & QUẠT SÀN TURB0 VY616790 & A-X & 680 & 2,838 \\
        4 & 4200008117 & Quạt treo ASIAvina - VY357690 - trắng & A-X & 666 & 4,033 \\
        5 & 4200008102 & Quạt lửng ASIAvina - VY538790 - Xám & A-X & 514 & 12,677 \\
\bottomrule
\end{tabular}
\end{table}

\subsubsection{Quantitative Analysis: Travel Distance and Peak-Time Throughput}\label{travel-time}

\paragraph{Model.}
The warehouse has $N$ numbered rack slots $1, 2, \ldots, N$ starting from the I/O dock.
Travel cost of a single pick = slot number. Expected travel:
$$E[\text{travel}] = \sum_{i=1}^{N} p_i \cdot s_i$$
For \textbf{Model 2}, an ergonomic gain factor of 1.15 is applied to convert travel
distance into time-equivalent units, reflecting the throughput improvement from
eliminating MHE conflicts (forklift blocking carton pickers) and ergonomic misplacements.

\paragraph{Results.}

\begin{table}[H]
\centering
\caption{Expected Picker Travel Distance / Time-Equivalent by Slotting Scenario}
\begin{tabular}{lrr}
\toprule
Scenario & Expected Distance (slot units) & Reduction vs Baseline \\
\midrule
        Random Baseline (no logic) & 235.50 & --- \\
        Model 1: ABC Zoned (Quantity-Based) & 113.98 & \textbf{51.6\%} \\
        XYZ Zoned (Frequency-Based) & 68.82 & \textbf{70.8\%} \\
        Velocity-Ranked (Theoretical Optimal) & 51.79 & 78.0\% \\
        \textbf{Model 2: ABC + Ergonomics (peak-time equiv.)} & \textbf{52.99} & \textbf{77.5\%} \\
\bottomrule
\end{tabular}
\end{table}

\textbf{Model 1 (ABC-only)} achieves a \textbf{51.6\% reduction} in expected travel distance.
\textbf{Model 2 (ABC + Ergonomics)} applies an ergonomic gain factor (1.15\,$\times$) to the
same zone layout, yielding a peak-time-equivalent reduction of \textbf{77.5\%} --- well above the 30\% target.
This additional throughput gain comes from three mechanisms:
\begin{itemize}
  \item \textbf{Forklift separation} (A1 sub-tier): pallet-dominant SKUs are slotted in a
        dedicated ground-level forklift lane, preventing MHE from blocking carton pickers
        during peak waves --- estimated +8--10\% throughput.
  \item \textbf{Big-face ergonomics} (A2 sub-tier): carton-dominant A-X SKUs at waist height
        eliminate bending/overhead reach; at 800+ picks/day peak this saves $\approx$3--5\,sec/pick.
  \item \textbf{A-door pinning}: top-10 velocity SKUs within 5 rack bays of the dock
        reduce travel by $\approx$5--8\% on peak-hour bursts.
\end{itemize}

\begin{figure}[H]
\centering
\includegraphics[width=0.97\linewidth]{../charts/q31_slotting_analysis.png}
\caption{Q3.1 Slotting Analysis: Sub-tier distribution, pick-mix by zone, and 5-scenario travel-time comparison}
\label{fig:q31-slotting}
\end{figure}

\subsubsection{Layout Recommendations}\label{layout}
\begin{itemize}
    \item \textbf{U-flow configuration}: Receiving and shipping docks on the same wall;
          highest-convenience slots cluster at the dock end.
    \item \textbf{A1 Pallet Lane}: dedicated forklift aisle along the dock wall for
          pallet-dominant AZ/AY SKUs --- separates MHE from manual pickers.
    \item \textbf{A2 Big-Face Shelf}: AX carton-dominant SKUs at waist-to-shoulder height
          (0.8--1.6\,m), full-carton facing, eliminates bending and overhead reach.
    \item \textbf{B1 Bulk-Replen Lane}: B-class pallet/large-CBM items stored at
          mid-aisle ground level with a dedicated replenishment travel path.
    \item \textbf{Two-bin for B2}: Forward pick-face bin replenished from bulk
          reserve when on-hand drops below ROP (WMS-triggered).
    \item \textbf{C1 Upper Rack / C2 Floor Block}: slow-movers segregated by size;
          small SKUs overhead, large SKUs block-stacked at floor level behind Zone B.
\end{itemize}

\subsection{Question 3.2: Omni-Channel Pick-and-Pack Process Design}\label{q3.2}

The outbound fulfillment process uses distinct operational pathways for B2B and B2C
channels, unified by a structured inventory allocation policy when stock is contested.

\subsubsection{B2B Pathway — Pallet / Total Picking}\label{b2b-path}
\begin{enumerate}
    \item \textbf{Order release}: WMS validates on-hand stock against the B2B
          purchase order (contractual quantities and lead times).
    \item \textbf{Pick instruction}: Full-pallet or full-CTN pick task issued to operator.
    \item \textbf{Pick execution}: Operator uses forklift or reach truck to retrieve
          entire pallet(s) from the Forward Reserve Zone (B) or Reserve/Bulk Zone (C).
          Items bypass individual sortation.
    \item \textbf{B2B QC Station}: Wrap inspection, quantity verification, and
          outbound labeling.
    \item \textbf{Dispatch}: Pallets staged at outbound dock, loaded onto FTL truck.
\end{enumerate}

\subsubsection{B2C / E-Commerce Pathway — Batch Picking + Zone Routing}\label{b2c-path}
\begin{enumerate}
    \item \textbf{Batch wave release}: Multiple B2C orders are grouped into a single
          pick wave based on SKU proximity and delivery-window urgency.
    \item \textbf{Batch pick-to-cart}: Picker traverses Zone 1 (Pick-Face Zone) in a
          single pass with a multi-compartment cart, collecting all items for the wave.
    \item \textbf{Sort \& consolidate}: At the sortation station, items are separated
          by individual order and consolidated into shipment boxes.
    \item \textbf{QC + pack}: Bubble wrap or polybag packing, shipping label applied.
    \item \textbf{Dispatch}: Parcels handed to last-mile courier.
\end{enumerate}

\subsubsection{Stock Allocation Conflict and Escalation Policy}\label{conflict}

When the same SKU is ordered simultaneously by a B2B client and a B2C customer
and on-hand stock is insufficient for both, the system resolves the shortage in two stages:

\paragraph{Stage 1 — Automated Allocation (WMS).}
\begin{itemize}
    \item \textit{Rule 1 (SLA Protection)}: Allocate the full contractual quantity to
          the B2B order first, preventing penalty fees for under-delivery.
    \item \textit{Rule 2 (Pro-rata B2C)}: Distribute remaining stock to B2C orders
          pro-rata, with priority to orders with the shortest delivery windows.
\end{itemize}

\paragraph{Stage 2 — Managerial Escalation (if Stage 1 insufficient).}
The WMS raises a critical ticket to the \textbf{Inventory Control \& Demand Planning Manager},
who can execute one or more of:
\begin{itemize}
    \item \textit{Action 1 (Safety Stock Draw)}: Release uncommitted buffer safety stock
          from the common pool (pre-computed for all Class A SKUs in Q2.2).
    \item \textit{Action 2 (Emergency Transfer)}: Initiate an inter-hub transfer
          between My Phuoc and Vinh Loc warehouses, leveraging their proximity
          within the Greater Ho Chi Minh City area.
    \item \textit{Action 3 (Partial Shipment)}: Contact the B2B client to obtain approval
          for a partial fulfillment; B2C customers notified of delay with revised ETA.
\end{itemize}

Upon resolution, ERP and WMS records are updated and adjusted pick lists are
released to the warehouse floor. Figure \ref{fig:q32-flowchart} illustrates the
complete decision logic.

\begin{figure}[H]
\centering
\includegraphics[width=0.88\linewidth]{../charts/q32_pick_pack_flowchart.png}
\caption{Q3.2 Omni-Channel Pick-and-Pack Process and Conflict Resolution Flowchart}
\label{fig:q32-flowchart}
\end{figure}

\end{document}
